% AI-produced by GPT-6 Astra (OpenAI, ChatGPT Pro mode).
% Pending audit; uses the separately audited Galvin RP3 component theorem.
\subsection{A formal obstruction on all ordinary three-dimensional concordance components}

\begin{saresult}{A.PENDING2}{Proved under stated hypotheses}{A natural Casson--Sullivan quotient of ordinary TOP concordance components}
For every closed, connected real PL three-manifold $B$ there is a
natural group homomorphism
\[
 c_B:\pi_0P_{\mathrm{TOP}}(B)\longrightarrow H^2(B;\mathbb Z/2).
                                                        \tag{C3.1}
\]
It vanishes on the image of $\pi_0P_{\mathrm{PL}}(B)$ and on the image
of $\pi_0P_{\mathrm{def}}(B)$ for every fixed o-minimal expansion of
$\mathbb R$. Naturality means conjugation by a PL homeomorphism of
bases, with the corresponding pullback on cohomology. If $B$ is
oriented and $\pi_1B$ is finite cyclic, then $c_B$ is surjective.
In particular, when that cyclic group has even order,
\[
 \pi_0P_{\mathrm{TOP}}(B)/\ker c_B\cong\mathbb Z/2,
 \qquad
 \operatorname{im}\bigl(\pi_0P_{\mathrm{def}}(B)\bigr)
       \subseteq\ker c_B.
                                                        \tag{C3.2}
\]
Only this containment is asserted. The kernel is not identified with
the PL or definable image, and no normality of that image in the whole
TOP component group is claimed.
\end{saresult}
\begin{proof}
\textbf{The formal relative class when the boundary map is PL.}
Let $X=B\times I$ with its product PL structure. Write $R_B$ for the
group of TOP concordances $h$ whose top restriction is a PL
homeomorphism; they already fix the bottom pointwise. For $h\in R_B$
the mapping cylinder $M_h$ has a prescribed PL reduction of its stable
tangent microbundle on the full boundary: the two copies of $X$ have
their given PL structures, and the side is the mapping cylinder of the
PL boundary map $1\amalg h|_{B\times1}$. Define
\[
 k_B(h)=\sigma^{-1}k(M_h,\partial M_h)
             \in H^3(X,\partial X;\mathbb Z/2).
                                                        \tag{C3.3}
\]
This extends the boundary-identity class (G.5) of the preceding
theorem. It is still a formal PL obstruction on a five-dimensional
mapping cylinder. It invokes no classification or smoothing of actual
four-dimensional isotopies.

We record its formal properties with their justification. The stable
tangent PL reduction of $X$, pulled back by $h$ and compared through
the tangent microbundle isomorphism of $h$, differs from the original
reduction by a relative $\mathrm{TOP}/\mathrm{PL}$ class. On the
boundary the PL map $h|_{\partial X}$ gives the specified identification
of these reductions. One may first use the product-collar compression
of A.136 to make $h$ a product of its boundary map on a common collar.
The normal line then has its literal product PL identification. The
same compression at a fixed width is multiplicative: on the smaller
cylinder it is conjugation by one embedding, and on the removed collars
it is multiplication of the boundary maps. Thus the composition rule
below can be checked with common product collars for any finite
collection, and the family version uses one width over the compact
parameter space. Since stable $\mathrm{TOP}/\mathrm{PL}$ is
$K(\mathbb Z/2,3)$
\cite[Essay V, Theorem~5.5(I), printed p.~251]{KirbySiebenmann1977},
the relative difference lies in the group in (C3.3). There is no
orientation local system here: the given PL reduction supplies a
basepoint in each fiber, and the pointed self-equivalence space of
$K(\mathbb Z/2,3)$ is weakly contractible. Its components are
$\operatorname{Aut}(\mathbb Z/2)=1$, and its positive homotopy groups
are the reduced groups
$\widetilde H^{3-i}(K(\mathbb Z/2,3);\mathbb Z/2)=0$ for $i>0$.
Thus the pointed fiber bundle is untwisted, just as in A.136.
The interval obstruction is precisely the relative universal
class defining that equation. Concatenating two comparison homotopies
gives the difference-of-lifts rule
\[
             k_B(uv)=k_B(v)+v^*k_B(u),\qquad u,v\in R_B.
                                                        \tag{C3.4}
\]
This is also the mapping-cylinder stacking proof of Galvin's
Proposition~2.23, printed pp.~2508--2509, with PL reductions in place
of smooth reductions. The boundary identifications in this argument
are PL and compose, which is the reason it applies to all of $R_B$.
The same relative lift construction proves invariance along TOP
isotopies with fixed PL boundary map. It also proves invariance when
the boundary varies through a jointly PL family, since that family
supplies the PL boundary reduction over the additional parameter.
Finally,
\[
                 h\text{ PL}\quad\Longrightarrow\quad k_B(h)=0,
                                                        \tag{C3.5}
\]
by the actual PL structure on its mapping cylinder.

Every bottom-fixed concordance acts trivially on this relative
cohomology group. Its top map $a$ is homotopic to the identity by the
map $(x,t)\mapsto\operatorname{pr}_B h(x,t)$. In the long exact
sequence of $(X,\partial X)$, restriction is diagonal in both degrees
two and three. Consequently the connecting map identifies
\[
 H^3(X,\partial X;\mathbb Z/2)
  \cong
 \operatorname{coker}\bigl(
 H^2(B;\mathbb Z/2)\xrightarrow{\Delta}
 H^2(B;\mathbb Z/2)\oplus H^2(B;\mathbb Z/2)\bigr).
                                                        \tag{C3.6}
\]
The boundary action is the identity on the bottom and $a^*=1$ on the
top, so it is the identity on the cokernel. Thus (C3.4) is additive
on $R_B$. This proof uses mod-two coefficients and does not require
orientability of $B$.

\textbf{Vanishing of every endpoint-trace ambiguity.}
If $\delta_s$ is any TOP path from the identity to a PL map of $B$,
let
\[
        T_\delta(x,s)=(\delta_s(x),s)\in R_B.
                                                        \tag{C3.7}
\]
The full ordinary automorphism comparison (A.134.2) replaces
$\delta$ by a jointly PL path with the same endpoints, through a
homotopy of paths fixed at those endpoints. Explicitly, its
$\pi_0$ injectivity supplies a PL path between the endpoints, and its
$\pi_1$ surjectivity corrects that path by a PL loop to obtain the
prescribed relative TOP path class. The resulting trace isotopy fixes
the boundary and ends at a PL trace. Therefore
\[
                         k_B(T_\delta)=0.
                                                        \tag{C3.8}
\]
Only an ordinary path comparison in dimension three has entered.

\textbf{Normalization, its independence, and path invariance.}
For an arbitrary $h\in P_{\mathrm{TOP}}(B)$, let $f$ be its top map.
Surjectivity on components of (A.134.2) supplies a PL map $a$ and a
TOP path $\gamma_s:f\to a$. Set
\[
 E_\gamma(x,s)=(\gamma_s f^{-1}(x),s),\qquad
                       h^\gamma=E_\gamma h\in R_B.
                                                        \tag{C3.9}
\]
The family with $\gamma_s$ replaced by $\gamma_{ts}$ is a path from
the identity to $E_\gamma$ in the bottom-fixed group, so $h^\gamma$
is in the same ordinary free-top component as $h$.

For another choice $\gamma':f\to a'$, direct multiplication gives
\[
 h^{\gamma'}(h^\gamma)^{-1}=E_{\gamma'}E_\gamma^{-1}
                  =T_\delta,
 \qquad \delta_s=\gamma'_s\gamma_s^{-1}.
                                                        \tag{C3.10}
\]
Here $\delta_0=1$ and $\delta_1=a'a^{-1}$ is PL. Equations
(C3.4), (C3.6), and (C3.8) imply
$k_B(h^{\gamma'})=k_B(h^\gamma)$. Thus this value is independent of
every normalization choice.

To prove constancy on components, it suffices first to replace the
initial map by a normalization. Let $H_t$ be a TOP path with
$H_0=h'\in R_B$ and $H_1=k$. Its top path $f_t$ starts at the PL map
$a=f_0$. Define
\[
 T_t(x,s)=(f_{ts}a^{-1}(x),s),\qquad Q_t=T_t^{-1}H_t.
                                                        \tag{C3.11}
\]
Every $Q_t$ has the fixed PL top map $a$ and fixes the bottom. Hence
$k_B(Q_t)$ is constant. If a common collar is required for the
relative tangent-boundary convention, apply the product-collar
compression of A.136 to this compact family; its top map remains
$a$ throughout. At $t=1$ the map $Q_1$ is precisely a normalization
of $k$ obtained from the path
\[
                  \gamma_s=a f_s^{-1}f_1:f_1\longrightarrow a.
                                                        \tag{C3.12}
\]
It follows that the common normalized value is constant along $H$.
For a path with an unnormalized initial map, first concatenate the
reverse of its normalization path with the given path. This proves
that the value defines a function on $\pi_0P_{\mathrm{TOP}}(B)$.

Compose it with the canonical interval-suspension identification
$H^3(X,\partial X;\mathbb Z/2)\cong H^2(B;\mathbb Z/2)$ to obtain
$c_B$. If $h$ and $k$ are arbitrary, choose their normalized
representatives $h'$ and $k'$. The product $h'k'$ is a normalized
representative of the component of $hk$, since the two normalization
paths can be multiplied. Additivity on $R_B$ now proves
\[
                        c_B([hk])=c_B([h])+c_B([k]).
                                                        \tag{C3.13}
\]

For a PL concordance choose the constant endpoint path; its value is
zero by (C3.5). For a definable concordance over any fixed o-minimal
expansion of the real field, A.113 supplies a definable isotopy to a
PL concordance. It is also a TOP isotopy, so its value is zero.
Conjugation by a PL homeomorphism of bases carries every mapping
cylinder, boundary reduction, and trace in the construction to the
corresponding object for the other base. Naturality of the relative
universal obstruction proves the stated naturality of $c_B$.

Finally, for closed oriented $B$ with finite cyclic fundamental group,
the relative realization proved in (G.8)--(G.14) realizes every class
in $H^3(X,\partial X;\mathbb Z/2)$ by a map fixed on both ends.
It is already normalized, so $c_B$ is surjective. For even order this
group is $\mathbb Z/2$ by mod-two Poincare duality and abelianization,
giving (C3.2). No step identifies all zero-class concordances with PL
maps; in particular the quotient by the PL image has not been used
as a group.
\end{proof}
